\documentclass[11pt]{article}
\usepackage[a4paper,margin=1in]{geometry}
\usepackage{amsmath,amssymb,amsthm,booktabs}
\usepackage[T1]{fontenc}
\usepackage{lmodern}
\usepackage[hidelinks]{hyperref}
\setlength{\emergencystretch}{3em}

\theoremstyle{plain}
\newtheorem{theorem}{Theorem}
\newtheorem{lemma}[theorem]{Lemma}
\newtheorem{proposition}[theorem]{Proposition}
\newtheorem{corollary}[theorem]{Corollary}
\theoremstyle{remark}
\newtheorem{remark}[theorem]{Remark}

\title{A Counterexample to Conjecture 00000000045}
\author{%
  Feiyu\thanks{Independent researcher.}
  \and
  Claude (Opus 5)\thanks{Anthropic.}%
}
\date{2026-09-12}

\begin{document}
\maketitle

\begin{abstract}
Conjecture 00000000045 asserts that the prime-sum graph $G_n$ has a Hamilton
cycle for all sufficiently large $n$. It does not: $G_n$ has a Hamilton cycle
only when $n$ is even, so every odd $n\ge3$ is a counterexample and no
threshold works. The refutation is formalized in Lean~4 against Mathlib.
\end{abstract}

\section{Statement}

For $n\ge1$ let $G_n$ be the graph with vertex set $\{1,\ldots,n\}$ in which
$i\sim j$ if and only if $i\ne j$ and $i+j$ is prime. Conjecture 00000000045
asserts that $G_n$ contains a Hamilton cycle for all sufficiently large $n$.

A Hamilton cycle of $G_n$ is a cyclic arrangement
\[
  v_0,v_1,\ldots,v_{n-1}
\]
of all $n$ vertices --- that is, $i\mapsto v_i$ is a bijection from
$\mathbb Z/n\mathbb Z$ onto $\{1,\ldots,n\}$ --- such that $v_i\sim v_{i+1}$ for
every $i$, indices read modulo $n$. We use this description throughout; it is
the definition.

\section{Parity obstruction}

\begin{lemma}\label{lem:odd}
Let $n\ge3$ and let $v_0,\ldots,v_{n-1}$ be a Hamilton cycle of $G_n$. Then
$v_i+v_{i+1}$ is odd for every $i$.
\end{lemma}

\begin{proof}
Since $n\ge3$, the positions $i$ and $i+1$ are distinct modulo $n$, so
$v_i\ne v_{i+1}$. Both are at least $1$, hence $v_i+v_{i+1}\ge1+2=3$. By
hypothesis $v_i+v_{i+1}$ is prime, and a prime greater than $2$ is odd.
\end{proof}

\begin{theorem}\label{thm:even}
If $G_n$ has a Hamilton cycle and $n\ge3$, then $n$ is even.
\end{theorem}

\begin{proof}
Sum the edge labels around the cycle. Since $i\mapsto i+1$ is a permutation of
$\mathbb Z/n\mathbb Z$, we have $\sum_i v_{i+1}=\sum_i v_i$, and therefore
\begin{equation}\label{eq:double}
  \sum_{i\in\mathbb Z/n\mathbb Z}\bigl(v_i+v_{i+1}\bigr)
  \;=\;2\sum_{i\in\mathbb Z/n\mathbb Z}v_i ,
\end{equation}
which is even. On the other hand the left-hand side is a sum of $n$ terms, each
odd by Lemma~\ref{lem:odd}, so it is congruent to $n$ modulo $2$. Hence
$n\equiv0\pmod 2$.
\end{proof}

\begin{corollary}
Conjecture 00000000045 is false.
\end{corollary}

\begin{proof}
Given any threshold $N$, the integer $n=2N+3$ is odd and satisfies $n\ge N$ and
$n\ge3$. By Theorem~\ref{thm:even}, $G_n$ has no Hamilton cycle. So no
threshold has the asserted property.
\end{proof}

\begin{remark}[the same obstruction, phrased as bipartiteness]
If $i$ and $j$ have the same parity and $i\ne j$, then $i+j$ is even and at
least $1+3=4$, hence not prime. So every edge of $G_n$ joins an odd vertex to
an even vertex: $G_n$ is bipartite, with sides
\[
  O=\{\text{odd }i\le n\},\qquad E=\{\text{even }i\le n\},
  \qquad |O|=\lceil n/2\rceil,\quad |E|=\lfloor n/2\rfloor .
\]
A Hamilton cycle in a bipartite graph alternates between the sides, so it
requires $|O|=|E|$, which fails for every odd $n$. Theorem~\ref{thm:even} is
this argument in a form convenient to formalize: the double count
\eqref{eq:double} replaces the appeal to alternation.
\end{remark}

\section{What survives}

The obstruction is exactly parity, and it says nothing for even $n$. A direct
search finds a Hamilton cycle for every even $n$ in the range checked:

\begin{center}
\begin{tabular}{lrrrrrrrrr}
\toprule
$n$ & 3 & 4 & 5 & 6 & 7 & 8 & 9 & 10 & 11\\
\midrule
Hamilton cycle & no & yes & no & yes & no & yes & no & yes & no\\
\bottomrule
\end{tabular}

\medskip

\begin{tabular}{lrrrrrrrrr}
\toprule
$n$ & 12 & 13 & 14 & 15 & 16 & 17 & 18 & 19 & 20\\
\midrule
Hamilton cycle & yes & no & yes & no & yes & no & yes & no & yes\\
\bottomrule
\end{tabular}
\end{center}

So the natural repair is to restrict the conjecture to even $n$:

\begin{quote}
For every sufficiently large even $n$, $G_n$ has a Hamilton cycle.
\end{quote}

We do not address that statement, which is a well-studied question and, to our
knowledge, open; it is closely tied to the existence of ``prime circles'' and
is supported by extensive computation. The table above is computational
evidence only and is not part of the Lean development.

\section{Formal verification}

The accompanying Lean~4 project formalizes the refutation against Mathlib. The
conjecture is stated in the formalization, so the final theorem is a
formalization of its negation rather than of facts that merely entail it.

\begin{itemize}
  \item \texttt{cycSucc} --- the cyclic successor on positions \texttt{Fin n},
    with \texttt{cycSucc\_injective}, \texttt{cycSucc\_bijective} and
    \texttt{cycSucc\_ne} (a position is never its own successor once $n\ge3$).
  \item \texttt{HasPrimeHamiltonCycle} --- a Hamilton cycle of $G_n$, as the
    cyclic arrangement described above: a permutation $\sigma$ of the positions,
    the vertex at position $i$ being $\sigma(i)+1$, with
    $\texttt{Nat.Prime}\bigl((\sigma i+1)+(\sigma(\texttt{cycSucc } i)+1)\bigr)$
    for every $i$.
  \item \texttt{ConjectureHolds} --- the conjecture.
  \item \texttt{even\_of\_hasPrimeHamiltonCycle} --- Theorem~\ref{thm:even}.
    The reindexing \eqref{eq:double} is Mathlib's
    \texttt{Fintype.sum\_bijective} applied to \texttt{cycSucc}.
  \item \texttt{conjecture\_00000000045\_false} ---
    $\lnot$\,\texttt{ConjectureHolds}, by instantiating at $n=2N+3$.
\end{itemize}

Primality is Mathlib's \texttt{Nat.Prime} throughout. The project builds with
\texttt{lake build} on \texttt{leanprover/lean4:v4.33.1} with Mathlib
\texttt{v4.33.1}. It contains no \texttt{sorry}, no \texttt{admit} and no
\texttt{native\_decide}, and introduces no axioms of its own:
\texttt{\#print axioms conjecture\_00000000045\_false} reports exactly the
three standard axioms \texttt{propext}, \texttt{Classical.choice} and
\texttt{Quot.sound}.

\end{document}
